\setlength{\parskip}{0pt}%
\lstdefinestyle{artefactopm}{
  basicstyle=\ttfamily\fontsize{8}{9}\selectfont, columns=fixed, basewidth=0.5em, numbers=none,
  showstringspaces=false, upquote=true, breaklines=true, breakatwhitespace=true,
  postbreak=\mbox{\textcolor{gray}{$\hookrightarrow$}\space},
  aboveskip=0pt, belowskip=0pt, xleftmargin=0pt, xrightmargin=0pt, frame=none,
  escapeinside={(*}{*)},
  literate={á}{{\'a}}1 {é}{{\'e}}1 {í}{{\'\i}}1 {ó}{{\'o}}1 {ú}{{\'u}}1 {ñ}{{\~n}}1
           {Á}{{\'A}}1 {É}{{\'E}}1 {Í}{{\'I}}1 {Ó}{{\'O}}1 {Ú}{{\'U}}1
           {·}{{\textperiodcentered}}1 {—}{{\textemdash}}1 {×}{{$\times$}}1,
}%
\mdfdefinestyle{artefactopm}{
  linewidth=0.5pt, linecolor=gray!70, leftmargin=0pt, rightmargin=0pt,
  innerleftmargin=5pt, innerrightmargin=5pt, innertopmargin=2pt, innerbottommargin=3pt,
  skipabove=0pt, skipbelow=3.5pt,
  frametitlerule=true, frametitlerulewidth=0.5pt, frametitlerulecolor=gray!70,
  frametitlebackgroundcolor=gray!12, frametitlefont=\normalfont\footnotesize,
  frametitleaboveskip=2pt, frametitlebelowskip=1.5pt, innerlinewidth=0pt,
}%
\providecommand*{\artcorte}{\textcolor{gray}{\textnormal{[\ldots]}}}%
\providecommand*{\artok}{\tikz[baseline=-0.25ex,x=1ex,y=1ex]\draw[sem_verde,line width=0.8pt,line cap=round,line join=round](0,0.45)--(0.35,0.1)--(0.95,0.9);}%
\providecommand*{\arttitulo}[3]{\textbf{#1\enspace #2}\hfill\texttt{03\_datos/\allowbreak apis/\allowbreak postmark/\allowbreak #3}}%
\begin{mdframed}[style=artefactopm,frametitle={\arttitulo{1}{Datasheet}{datasheet.yaml}}]
\begin{lstlisting}[style=datasheetyaml, style=artefactopm]
capacity:
  (*\artcorte*)
  quota_email_included_basic:
    value: 10000                                            # (*\semambar*)
    (*\artcorte*)
    overage_cost:
      price: 0.0018                                         # (*\semverde*)
      value: 1
(*\artcorte*)
plans:
  (*\artcorte*)
  basic:
    price: 15.00                                            # (*\semverde*)
    (*\artcorte*)
    quota: quota_email_included_basic
    rate: rate_request_rest
\end{lstlisting}
\end{mdframed}
\begin{mdframed}[style=artefactopm,frametitle={\arttitulo{2}{Registro de fuentes}{fuentes.md}}]
\begin{lstlisting}[style=artefactopm]
## [F1] Página de pricing

- **URL:** https://postmarkapp.com/pricing
- **Fecha de acceso:** 2026-08-03 (re-verificada; primera visita 2026-08-02)
- **Citas literales:**
  (*\artcorte*)
  - Basic — `"$15.00/mo"`, `"Starting at 10,000 emails/month"`,
    `"Extra emails @ $1.80 / 1,000"`
\end{lstlisting}
\end{mdframed}
\begin{mdframed}[style=artefactopm,frametitle={\arttitulo{3}{Registro de elementos no soportados}{fricciones.md}}]
\begin{lstlisting}[style=artefactopm]
## F-02 · Cuotas por escalón, no por valor único — y con overage por escalón

**Qué pasa.** Los tres planes de pago dicen *"Starting at 10,000 emails/month"* [F1].
El panel [F11] revela la estructura completa: una escalera de volumen (10k/$15 ·
50k/$55 · 125k/$115 · 300k/$245 · 700k/$455 · 1,5M/$775) en la que **el propio
overage depende del escalón** ($1.80/1.000 hasta 125k, $1.70/1.000 desde 300k), y
(*\artcorte*)
**Qué hace v0.3.** Una cuota tiene un `value` escalar y un `overage_cost` único. Ni
la escalera ni el overage-por-escalón caben.
(*\artcorte*)
**Propuesta.** Permitir que una cuota declare escalones (`tiers`) con precio y
overage propios por escalón, más un escalón `custom` para el tramo negociado.
\end{lstlisting}
\end{mdframed}
\begin{mdframed}[style=artefactopm,frametitle={\arttitulo{4}{Acta de revisión}{REVISION.md}}]
\begin{lstlisting}[style=artefactopm]
### Q9 · Proyección del producto tier × volumen en `plans` (F-08) — (*\artok*) DECIDIDA (*\artcorte*)
El panel [F11] revela que Basic/Pro/Platform son **tiers funcionales** y el volumen
va aparte, en escalera. La datasheet mantiene la proyección "3 planes = 3 tiers, con
la cuota del escalón mínimo" y documenta la escalera en provenance/F-02. Alternativa
descartada: enumerar las ~18 combinaciones tier×escalón como planes. (*\artcorte*)
\end{lstlisting}
\end{mdframed}
\begin{mdframed}[style=artefactopm,frametitle={\arttitulo{5}{Matriz de \textit{features} por plan}{features.md}}]
\begin{lstlisting}[style=artefactopm]
## Features por plan

| Feature | Free | Basic | Pro | Platform | Fuente · escalón |
|---|---|---|---|---|---|
(*\artcorte*)
| Usuarios | — | `"Up to 4 users"` | `"Up to 6 users"` | `"Unlimited users"` | [F11] panel |
\end{lstlisting}
\end{mdframed}
